\begin{abstract}
We report a scoped decomposition of spatial question answering over real handheld room
captures. Rather than a single accuracy number, we measure each pipeline stage against a
scope discipline separating what a system \emph{could} express from what a deployable
path can \emph{reach}, and find the two far apart.

Replacing only the labeler's input images --- real capture crops instead of point-splat
renders, with weights, vocabulary, threshold, evaluator, delivered partitions and
matching all fixed --- raises matched-instance top-1 from 1/21 to 12/21
and top-3 from 4/21 to 18/21 across three scenes with no tuning between
them. Paired per instance the change is one-directional. This is a
component evaluation on an oracle-selected denominator, not end-to-end performance.

On a relation benchmark over two rooms, the \emph{same stored relations} answer 7 of 10
scored items when identity is supplied by a human, 0 of 10 from learned labels,
and 2 of 10 through an oracle-free grounding bridge --- a paired
substitution sizing a bound, not a system improvement. A replay reading only serialized
edges agrees with recomputed geometry on 12 of 12 items: no additional loss was measured
during serialization relative to recomputation under the same AABB convention. Relation
availability and serialization consistency were observed on those items; semantic
correctness of the stored relations was not independently established. The stored
content was not reached by any of the three identity paths we instantiated; the binding
stage is identity.

On a blinded one-scene pilot room, unused during development, direct multiview RGB
answers 5 of 10 with \textbf{zero wrong answers}, declines all three non-co-visible
items, and fails both predeclared gates.

The contribution is the decomposition and its scope discipline, with negative results
committed as first-class artifacts --- not a solved room-understanding system.
\end{abstract}
